% GENERATED FILE. Edit canonical lessons, then run npm run generate:books.
% canonical-source-hash: fnv1a64:0e3cc37f26026ecf
% canonical-lessons: ZH-C492-yuwen, ZH-C492-shuxue, ZH-C492-dili, ZH-C492-yingyu, ZH-C492-zhuanye

\chapter{School Subjects}
\label{ch:zh-school-subjects}
\hlchaptermodality{492}

\begin{chapteropening}
\textbf{By the end of this chapter:} \emph{I can name school subjects, say which one I like best and talk about my major.}

Complete the last lesson of chapter 492: I can name school subjects, say which one I like best and talk about my major.
\end{chapteropening}

\section[\texorpdfstring{\zh{语文} (yǔwén)}{语文 (yǔwén)}]{\zh{语文} (yǔwén) — Chinese (as a school subject)}

\label{lesson:ZH-C492-yuwen}



\begin{quote}
Before the new one: say the Mandarin for the world of work, then the Mandarin for someone under one at work.
\end{quote}

\subsection*{You'll want to know: \zh{语文}}
\textbf{\zh{语文}} — \emph{yǔwén} — \textquotedblleft{}Chinese (as a school subject)\textquotedblright{}.

\begin{quote}
\textbf{\zh{我最喜欢语文课}} \emph{wǒ zuì xǐhuan yǔwén kè} — Chinese is my favourite lesson.
\end{quote}

\begin{cousinweb}
\textbf{\zh{语文}} (\emph{yǔwén}) is built from \zh{语} 'speech' and \zh{文} 'writing', both inherited from Middle Chinese, \emph{ngjoX} and \emph{mjun} (Guangyun, in Baxter's transcription), which glosses the first 'to discuss' and the second 'a written piece'. The name of the school subject is modern and not given there.
\end{cousinweb}

\begin{grammarlens}[title={What you've built}]
One more word for saying how it was.
\end{grammarlens}

\subsection*{Your turn}
\begin{itemize}
\raggedright
  \item \emph{Say it:} \emph{yǔwén}
  \item \emph{Say it:} \emph{yǔwén}, once more
  \item \emph{Say it:} \emph{wǒ zuì xǐhuan yǔwén kè} — Chinese is my favourite lesson.
  \item \emph{Recall:} say the Mandarin for the world of work, then the Mandarin for someone under one at work, then say \emph{yǔwén} again
\end{itemize}

\begin{tcolorbox}[breakable,colback=teal!4,colframe=teal!35!black,title={Before you move on}]
What does \zh{语文} mean? (\textquotedblleft{}Chinese (as a school subject)\textquotedblright{}.) Say it once more.
\end{tcolorbox}
\section[shùxué]{shùxué — maths; mathematics}

\label{lesson:ZH-C492-shuxue}



\begin{quote}
Before the new one: say the Mandarin for someone under one at work, then the Mandarin for Chinese (as a school subject).
\end{quote}

\subsection*{You'll want to know: shùxué}
\emph{shùxué} — \textquotedblleft{}maths; mathematics\textquotedblright{}.

\begin{quote}
\emph{wǒ de shùxué bǐ yǔwén hǎo} — I'm better at maths than at Chinese.
\end{quote}

Its characters come later; for now it is a word for the ear.

\begin{cousinweb}
\emph{shùxué} is built from \emph{shù} 'a number' and \zh{学} 'to learn', both inherited from Middle Chinese, \emph{srjuH} and \emph{haewk} (Guangyun, in Baxter's transcription), which glosses the first 'reckoning' and names the nine arts of reckoning of an old book of rites.
\end{cousinweb}

\begin{grammarlens}[title={What you've built}]
One more word for saying how it was.
\end{grammarlens}

\subsection*{Your turn}
\begin{itemize}
\raggedright
  \item \emph{Say it:} \emph{shùxué}
  \item \emph{Say it:} \emph{shùxué}, once more
  \item \emph{Say it:} \emph{wǒ de shùxué bǐ yǔwén hǎo} — I'm better at maths than at Chinese.
  \item \emph{Recall:} say the Mandarin for someone under one at work, then the Mandarin for Chinese (as a school subject), then say \emph{shùxué} again
\end{itemize}

\begin{tcolorbox}[breakable,colback=teal!4,colframe=teal!35!black,title={Before you move on}]
What does shùxué mean? (\textquotedblleft{}Maths; mathematics\textquotedblright{}.) Say it once more.
\end{tcolorbox}
\section[\texorpdfstring{\zh{地理} (dìlǐ)}{地理 (dìlǐ)}]{\zh{地理} (dìlǐ) — geography}

\label{lesson:ZH-C492-dili}



\begin{quote}
Before the new one: say the Mandarin for Chinese (as a school subject), then the Mandarin for maths.
\end{quote}

\subsection*{You'll want to know: \zh{地理}}
\textbf{\zh{地理}} — \emph{dìlǐ} — \textquotedblleft{}geography\textquotedblright{}.

\begin{quote}
\emph{wǒmen de dìlǐ lǎoshī hěn yǒuyìsi} — Our geography teacher is really interesting.
\end{quote}

\begin{cousinweb}
\textbf{\zh{地理}} (\emph{dìlǐ}) is built from \zh{地} 'the land' and \zh{理} 'a pattern', both inherited from Middle Chinese, \emph{dijH} and \emph{liX} (Guangyun, in Baxter's transcription), which glosses the first 'the earth' and the second 'to work jade' and 'a pattern'.
\end{cousinweb}

\begin{grammarlens}[title={What you've built}]
One more word for saying how it was.
\end{grammarlens}

\subsection*{Your turn}
\begin{itemize}
\raggedright
  \item \emph{Say it:} \emph{dìlǐ}
  \item \emph{Say it:} \emph{dìlǐ}, once more
  \item \emph{Say it:} \emph{wǒmen de dìlǐ lǎoshī hěn yǒuyìsi} — Our geography teacher is really interesting.
  \item \emph{Recall:} say the Mandarin for Chinese (as a school subject), then the Mandarin for maths, then say \emph{dìlǐ} again
\end{itemize}

\begin{tcolorbox}[breakable,colback=teal!4,colframe=teal!35!black,title={Before you move on}]
What does \zh{地理} mean? (\textquotedblleft{}Geography\textquotedblright{}.) Say it once more.
\end{tcolorbox}
\section[Yīngyǔ]{Yīngyǔ — English (the language)}

\label{lesson:ZH-C492-yingyu}



\begin{quote}
Before the new one: say the Mandarin for maths, then the Mandarin for geography.
\end{quote}

\subsection*{You'll want to know: Yīngyǔ}
\emph{Yīngyǔ} — \textquotedblleft{}English (the language)\textquotedblright{}.

\begin{quote}
\emph{wǒ xué le liù nián Yīngyǔ} — I studied English for six years.
\end{quote}

Its characters come later; for now it is a word for the ear.

\begin{cousinweb}
\emph{Yīngyǔ} is built from \emph{Yīng}, short for the name of England, and \zh{语} 'speech', both inherited from Middle Chinese, \emph{'jaeng} and \emph{ngjoX} (Guangyun, in Baxter's transcription), which glosses the first 'a flower'; the first stands for the country only in modern use.
\end{cousinweb}

\begin{grammarlens}[title={What you've built}]
One more word for saying how it was.
\end{grammarlens}

\subsection*{Your turn}
\begin{itemize}
\raggedright
  \item \emph{Say it:} \emph{Yīngyǔ}
  \item \emph{Say it:} \emph{Yīngyǔ}, once more
  \item \emph{Say it:} \emph{wǒ xué le liù nián Yīngyǔ} — I studied English for six years.
  \item \emph{Recall:} say the Mandarin for maths, then the Mandarin for geography, then say \emph{Yīngyǔ} again
\end{itemize}

\begin{tcolorbox}[breakable,colback=teal!4,colframe=teal!35!black,title={Before you move on}]
What does Yīngyǔ mean? (\textquotedblleft{}English (the language)\textquotedblright{}.) Say it once more.
\end{tcolorbox}
\section[zhuānyè]{zhuānyè — major; field of study}

\label{lesson:ZH-C492-zhuanye}



\begin{quote}
Before the new one: say the Mandarin for geography, then the Mandarin for English (the language).
\end{quote}

\subsection*{You'll want to know: zhuānyè}
\emph{zhuānyè} — \textquotedblleft{}major; field of study\textquotedblright{}.

\begin{quote}
\emph{nǐ de zhuānyè shì shénme} — What's your major?
\end{quote}

Its characters come later; for now it is a word for the ear.

\begin{cousinweb}
\emph{zhuānyè} is built from \emph{zhuān} 'single-minded' and \zh{业} 'a trade', both inherited from Middle Chinese, \emph{tsywen} and \emph{ngjaep} (Guangyun, in Baxter's transcription), which glosses the first 'alone' and 'to have to oneself'.
\end{cousinweb}

\begin{grammarlens}[title={What you've built}]
One more word for saying how it was.
\end{grammarlens}

\subsection*{Your turn}
\begin{itemize}
\raggedright
  \item \emph{Say it:} \emph{zhuānyè}
  \item \emph{Say it:} \emph{zhuānyè}, once more
  \item \emph{Say it:} \emph{nǐ de zhuānyè shì shénme} — What's your major?
  \item \emph{Recall:} say the Mandarin for geography, then the Mandarin for English (the language), then say \emph{zhuānyè} again
\end{itemize}

\begin{tcolorbox}[breakable,colback=teal!4,colframe=teal!35!black,title={Before you move on}]
What does zhuānyè mean? (\textquotedblleft{}Major; field of study\textquotedblright{}.) Say it once more.
\end{tcolorbox}
